% Natarajan 打散 = 标签网格中的完整四角矩形。
% 左：六模式族 V（同 lt6-email-label-grid）。试取 x1 的标签对{正常,广告}、
%     x2 的标签对{正常,广告}：角 (广告,正常) 缺失。V 中任意两行只共享一列，
%     故任何选法都缺角，d_N(V)=1。
% 右：V+ = V ∪ {(正常,诈骗)}（黄色点为新增）。x1 取{正常,广告}，x2 取{广告,诈骗}，
%     四角 (正常,广告)(正常,诈骗)(广告,广告)(广告,诈骗) 均在 V+ 中，故 {x1,x2} 被
%     Natarajan 打散，d_N(V+)=2。
\begingroup
\begin{tikzpicture}[figure picture,x=1mm,y=1mm]
 \def\cs{10.5}
 \newcommand{\ltsixgrid}[2]{% #1 = x offset, #2 = heading
  \foreach \r in {1,2,3}{\foreach \c in {1,2,3}{
    \draw[FigureGrid,line width=.5pt] ({#1+\cs*(\c-1)},{21-\cs*(\r-1)}) rectangle ++(\cs,\cs);}}
  \foreach \lab [count=\c] in {正常,广告,诈骗}{
    \node[anchor=south,font=\figurefont\scriptsize] at ({#1+\cs*(\c-.5)},32) {\lab};}
  \foreach \lab [count=\r] in {正常,广告,诈骗}{
    \node[anchor=east,font=\figurefont\scriptsize] at ({#1-1},{26.25-\cs*(\r-1)}) {\lab};}
  \node[figure heading] at ({#1+1.5*\cs},41) {#2};
 }
 % ---- left panel
 \ltsixgrid{9}{$\mathcal V$：任意两行只共一列}
 \foreach \r/\c in {1/1,1/2,2/2,2/3,3/3,3/1}{
   \fill[FigureBlue] ({9+\cs*(\c-.5)},{26.25-\cs*(\r-1)}) circle (1.7);}
 \foreach \r/\c in {1/3,3/2}{
   \draw[FigureMuted,line width=.6pt,fill=white] ({9+\cs*(\c-.5)},{26.25-\cs*(\r-1)}) circle (1.3);}
 \draw[FigureRed,line width=.85pt,dashed,rounded corners=1pt]
   ({9+\cs*.5},26.25) rectangle ({9+\cs*1.5},{26.25-\cs});
 \draw[FigureRed,line width=.85pt] ({9+\cs*.5-1.6},{26.25-\cs-1.6})--++(3.2,3.2)
   ({9+\cs*.5-1.6},{26.25-\cs+1.6})--++(3.2,-3.2);
 \node[figure note,anchor=north] at ({9+1.5*\cs},-1.5) {角$(\text{广告},\text{正常})$缺失};
 % ---- right panel
 \ltsixgrid{60}{$\mathcal V^{+}=\mathcal V\cup\{(\text{正常},\text{诈骗})\}$}
 \foreach \r/\c in {1/1,1/2,2/2,2/3,3/3,3/1}{
   \fill[FigureBlue] ({60+\cs*(\c-.5)},{26.25-\cs*(\r-1)}) circle (1.7);}
 \draw[FigureMuted,line width=.6pt,fill=white] ({60+\cs*1.5},{26.25-\cs*2}) circle (1.3);
 \draw[FigureMuted,line width=.6pt,fill=white] ({60+\cs*.5},{26.25-\cs}) circle (1.3);
 \fill[FigureYellow] ({60+\cs*2.5},26.25) circle (1.7);
 \draw[FigureBlue,line width=1pt,rounded corners=1pt]
   ({60+\cs*1.5},26.25) rectangle ({60+\cs*2.5},{26.25-\cs});
 \node[figure note,anchor=north] at ({60+1.5*\cs},-1.5) {$x_1$取正常或广告，$x_2$取广告或诈骗\\四种组合齐全};
\end{tikzpicture}
\endgroup
